% !TEX root = task-2-6-6.tex
% Editable task bank. Append new BankTask blocks at the end of the document.
\documentclass[12pt,a4paper]{article}
\usepackage{physics-textbook}
\usepackage{graphicx,needspace}
\usetikzlibrary{decorations.pathmorphing}
\newif\ifShowSolutions
\ifdefined\TutorEdition
  \ShowSolutionstrue
\else
  \ShowSolutionsfalse % Change to \ShowSolutionstrue to display solutions.
\fi
\newcounter{banktask}
\newsavebox{\BankTaskBox}
\newcommand{\BankTaskContent}[3]{%
  \noindent\textbf{\thebanktask.}\quad #2\par
  \ifShowSolutions
    \Needspace{50mm}% Keep the solution heading with its opening text or diagram.
    #3\par
    \ifdefined\TutorEdition
      \ifcsname TutorSource#1\endcsname
        \nopagebreak\smallskip{\small\csname TutorSource#1\endcsname\par}%
      \fi
    \fi
  \fi
}
\newcommand{\BankTask}[3]{%
  \refstepcounter{banktask}\par\medskip
  \sbox{\BankTaskBox}{\begin{minipage}{\linewidth}%
    \BankTaskContent{#1}{#2}{#3}\end{minipage}}%
  \ifdim\dimexpr\ht\BankTaskBox+\dp\BankTaskBox\relax>\dimexpr\textheight-20mm\relax
    \Needspace{35mm}\BankTaskContent{#1}{#2}{#3}%
  \else
    \Needspace{\dimexpr\ht\BankTaskBox+\dp\BankTaskBox\relax}%
    \noindent\usebox{\BankTaskBox}\par
  \fi
  \medskip
}
\begin{document}
\BookSubsection{2.6.6}{Изохорный процесс}


% Source sdamgia-9500: https://phys-ege.sdamgia.ru/problem?id=9500
\BankTask{sdamgia-9500}{
В закрытом сосуде с жесткими стенками содержится идеальный газ при температуре $27^\circ\text{C}$. Температуру газа повысили до $147^\circ\text{C}$. Во сколько раз изменилось давление этого газа? (До десятых долей.)
}{
\smallskip\noindent\textbf{Решение}\par
Идеальный газ в сосуде нагревается изохорно. Следовательно, согласно закону Шарля $( \frac{p}{T} =const ) ,$ давление газа должно увеличиться. Температура газа увеличивается в $\frac{T_2}{T_1} = \frac{147 + 273}{27 + 273} =1{,}4$ раза. А значит, и давление увеличится в $1{,}4$ раза.

Ответ: $1{,}4$.
}

% Source sdamgia-29027: https://phys-ege.sdamgia.ru/problem?id=29027
\BankTask{sdamgia-29027}{
В жестком сосуде находится некоторое постоянное количество идеального газа. Температуру газа увеличили в $3$ раза, а давление при этом увеличилось на $50\,\text{кПа}$. Определите начальное давление газа в сосуде.

(В кПа.)
}{
\smallskip\noindent\textbf{Решение}\par
Поскольку сосуд жесткий, то его объем не меняется, то есть процесс является изохорным. Для этого процесса выполняется закон Шарля $\frac{p_1}{p_2} = \frac{T_1}{T_2} ,$ причем $p_2=p_1 + \Delta p$ и $T_2=3T_1.$ Отсюда получаем $3p_1=p_1 + \Delta p.$ Значит:

$p_1= \frac{\Delta p}{2} = \frac{50}{2} =25\text{кПа}.$

Ответ: $25$.
}

% Source sdamgia-3297: https://phys-ege.sdamgia.ru/problem?id=3297
\BankTask{sdamgia-3297}{
При повышении температуры газа в запаянном сосуде давление газа увеличивается. Это изменение давления объясняется тем, что

$1$) увеличивается объем сосуда за счет нагревания его стенок

$2$) увеличивается энергия теплового движения молекул газа

$3$) увеличиваются размеры молекул газа при его нагревании

$4$) увеличивается энергия взаимодействия молекул газа друг с другом
}{
\smallskip\noindent\textbf{Решение}\par
Молекулы газа непрерывно ударяются о стенки сосуда, оказывая на них давление. Чем больше средняя кинетическая энергия теплового движения молекул, тем быстрее молекулы двигаются и тем больший импульс передают стенкам сосуда в единицу времени, оказывая тем самым большее давление. С другой стороны, температура является мерой энергии теплового движения молекул: $\overline{E}= \frac{3}{2} kT.$ Отсюда заключаем, что при повышении температуры газа в запаянном сосуде давление газа увеличивается вследствие увеличения энергии теплового движения молекул. 

Ответ: $2$.
}

% Source sdamgia-39946: https://phys-ege.sdamgia.ru/problem?id=39946
\BankTask{sdamgia-39946}{
В ходе изохорного процесса, происходящего с постоянной массой идеального газа, его давление возросло с $40\,\text{кПа}$ до $120\,\text{кПа}$. Во сколько раз увеличилась средняя кинетическая энергия теплового движения молекул газа в этом процессе?
}{
\smallskip\noindent\textbf{Решение}\par
При изохорном процессе выполняется закон Шарля:

$\frac{p_2}{p_1} = \frac{T_2}{T_1} = \frac{120}{40} =3.$

Таким образом, абсолютная температура газа увеличивается в $3$ раза. Средняя кинетическая энергия теплового движения молекул $\overline E_{\text{к}}= \frac{3}{2} kT.$ При увеличении температуры в $3$ раза средняя кинетическая энергия молекул увеличится в $3$ раза.

Ответ: $3$.
}

% Source sdamgia-43061: https://phys-ege.sdamgia.ru/problem?id=43061
\BankTask{sdamgia-43061}{
В закрытом сосуде постоянного объема некоторое количество теплоты было отведено от идеального газа. Как в результате этого изменились давление газа и его внутренняя энергия? 

$1$) увеличилась

$2$) уменьшилась

$3$) не изменилась

\begin{center}\begin{tabular}{|p{0.43\linewidth}|p{0.43\linewidth}|}\hline
Давление & Внутренняя энергия\\\hline
\strut & \strut\\\hline\end{tabular}\end{center}
}{
\smallskip\noindent\textbf{Решение}\par
Так как газ имеет постоянный объем, то работа газа равна нулю. По первому закону термодинамики отведенное количество теплоты привело к уменьшению внутренней энергии ($2$).

При уменьшении внутренней энергии газа его абсолютная температура изохорно уменьшилась, что привело к уменьшению давления газа($2$).

Ответ: $22$.
}

% Source sdamgia-7111: https://phys-ege.sdamgia.ru/problem?id=7111
\BankTask{sdamgia-7111}{
($N$ — число частиц, $p$ — давление, $V$ — объем, $T$ — абсолютная температура, $Q$ — количество теплоты)

\TaskMatchingColumns{Процессы}{А)  Изохорный процесс при $N=const$\par Б)  Изотермический процесс при $N=const$}{Формулы}{$1$)  $\frac{p}{T} = const$\par $2$)  $\frac{V}{T} = const$\par $3$)  $pV = const$\par $4$)  $Q = 0$}

\begin{center}\begin{tabular}{|c|c|}\hline
$A$ & Б\\\hline
\strut & \strut\\\hline\end{tabular}\end{center}
}{
\smallskip\noindent\textbf{Решение}\par
Уравнение состояния идеального газа:$pV=\nu RT$

Изохорный процесс означает, что объем газа остается постоянным, а значит, должно выполняться $\frac{p}{T} =const.$

Изотермический процесс означает, что температура газа остается постоянной, а значит, должно выполняться условие $pV=const.$

Ответ: $13$.
}

% Source sdamgia-11269: https://phys-ege.sdamgia.ru/problem?id=11269
\BankTask{sdamgia-11269}{
Один моль идеального одноатомного газа, находившегося при давлении $p_{1}$ в сосуде объемом $V_{1}$, изобарно нагревают от температуры $T_{1}$ до температуры $T_{2}$. Как изменятся приращение внутренней энергии газа и переданное газу количество теплоты, если нагревание этого газа осуществлять изохорно из того же начального состояния ($p_{1}$, $V_{1}$, $T_{1}$) до той же конечной температуры $T_{2}$?

$1$. Увеличится.

$2$. Уменьшится.

$3$. Не изменится.

\begin{center}\begin{tabular}{|p{0.43\linewidth}|p{0.43\linewidth}|}\hline
Приращение внутренней энергии газа & Переданное газу количество теплоты\\\hline
\strut & \strut\\\hline\end{tabular}\end{center}
}{
\smallskip\noindent\textbf{Решение}\par
Внутренняя энергия идеального газа связана с температурой соотношением: 

$U = \frac{3}{2} \nu R T.$ 

Таким образом, приращение внутренней энергии зависит только от изменения температуры газа. В обоих процессах температура изменяется на одинаковую величину, значит, приращение внутренней энергии не изменится.

Согласно первому началу термодинамики в изобарном процессе внешнее тепло идет на совершение газом работы и увеличение его внутренней энергии:

$Q=\Delta U + A.$

В изохорном процессе газ не совершает работы, значит, ему нужно передать меньшее количество теплоты для нагрева.

Ответ: $32$.
}

% Source sdamgia-3438: https://phys-ege.sdamgia.ru/problem?id=3438
\BankTask{sdamgia-3438}{
Идеальный газ изохорно нагревают так, что его температура изменяется на $\Delta T=240\text{К},$ а давление — в $1{,}6$ раза. Масса газа постоянна. Какова начальная температура газа по шкале Кельвина?
}{
\smallskip\noindent\textbf{Решение}\par
Изохорный процесс — это процесс, проводимый при постоянном объеме. Идеальный газ подчиняется уравнению Клапейрона — Менделеева. В начальном состоянии (до нагревания): $p_1V=\nu RT_1.$ После нагревания: $p_2V=\nu RT_2.$ Вычтя из второго уравнения первое и приняв во внимание, что $p_2=1{,}6p_1.$ получаем соотношение $0{,}6p_1V=\nu R\Delta T,$ откуда $\frac{p_1V}{\nu R} = \frac{\Delta T}{0{,}6} .$ Тогда начальная температура газа равна $T_1 = \frac{p_1V}{\nu R} = \frac{\Delta T}{0{,}6} = \frac{240\text{К}}{0{,}6} =400\text{К}.$

Ответ: $400\,\text{К}$.
}

% Source sdamgia-3890: https://phys-ege.sdamgia.ru/problem?id=3890
\BankTask{sdamgia-3890}{
Идеальный одноатомный газ, находящийся при температуре $T$, нагрели до температуры $2$ $T$, сообщив ему количество теплоты $10\,\text{Дж}$. В результате гaз совершил работу $5\,\text{Дж}$. Какое количество теплоты отдаст газ, если его после этого изохорически охладить до температуры $\frac{3T}{2} .$ (В джоулях, округлите до десятых.)
}{
\smallskip\noindent\textbf{Решение}\par
Согласно первому началу термодинамики, переданное газу тепло связано с изменением внутренней энергией и работой газа против внешних сил соотношением: $Q=A + \Delta U.$ Для идеального одноатомного газа изменение внутренней энергии определяется только изменением температуры: $\Delta U= \frac{3}{2} \nu R \Delta T .$ 

Таким образом, изменение внутренней энергии в первом процессе равно:

$\Delta U_1= \frac{3}{2} \nu R ( 2T - T ) = \frac{3}{2} \nu R T= Q_1 - A_1=10\text{Дж} - 5\text{Дж}=5 \text{Дж}$

Для второго процесса имеем (работа равна нулю, поскольку процесс идет изохорически):

$Q_2=\Delta U_2= \frac{3}{2} \nu R ( 3T/2 - 2T ) = - \frac{1}{2} \cdot \frac{3}{2} \nu R T= - \frac{1}{2} \cdot 5\text{Дж}= - 2{,}5 \text{Дж}$

Следовательно, при данном изохорическом охлаждении газ отдаст $2{,}5\,\text{Дж}$.

Ответ: $2{,}5\,\text{Дж}$.
}

% Source sdamgia-9753: https://phys-ege.sdamgia.ru/problem?id=9753
\BankTask{sdamgia-9753}{
Неизменное количество идеального одноатомного газа изохорически переводят из состояния $1$ в состояние $2$. Затем газ изобарически переводят из состояния $2$ в состояние $3$, и при этом газ совершает работу $40\,\text{Дж}$. Известно, что температура газа в процессе $2$-$3$ повышается на столько же, на сколько она повысилась в процессе $1$-$2$. Какое количество теплоты поглотил газ в процессе $1$-$2$?
}{
\smallskip\noindent\textbf{Решение}\par
В изобарном процессе совершается работа $A_{23}=\nu R\Delta T.$ 

В изохорном процессе работа газа равна нулю, и все тепло идет исключительно на увеличение внутренней энергии:

$Q_{12}=\Delta U_{12}= \frac{3}{2} \nu R\Delta T= \frac{3}{2} A_{23}= \frac{3}{2} \cdot 40 =60\text{Дж}.$

Ответ: $60\,\text{Дж}$.
}


% Source 349197: https://ege.fipi.ru/bank/questions.php?search=1&proj=BA1F39653304A5B041B656915DC36B38&qid=349197
\BankTask{349197}{
Идеальный газ в количестве \(5\ \text{моль}\) нагревается от \(t_1=127\,^{\circ}\text{С}\) до \(t_2=227\,^{\circ}\text{С}\) в сосуде постоянного объёма. Чему равно отношение давлений \(\frac{p_1}{p_2}\) идеального газа?
}{
\smallskip\noindent\textbf{Решение от GPT}\par
При постоянных объёме и количестве вещества \(p/T=\mathrm{const}\). Температуры: \(T_1=127+273=400\ \text{К}\), \(T_2=227+273=500\ \text{К}\). Тогда
\[\frac{p_1}{p_2}=\frac{T_1}{T_2}=\frac{400}{500}=0{,}8.\]
Ответ: \(0{,}8\).
}

% Source 67AF6A: https://ege.fipi.ru/bank/questions.php?search=1&proj=BA1F39653304A5B041B656915DC36B38&qid=67AF6A
\BankTask{67AF6A}{
С идеальным газом происходит изохорный процесс, в котором в результате увеличения абсолютной температуры газа в $2$ раза его давление возросло на $75\,\text{кПа}$. Масса газа постоянна. Каково было первоначальное давление газа (в $\text{кПа}$)?
}{
\smallskip\noindent\textbf{Решение от GPT}\par
При постоянных массе и объёме $p/T=\mathrm{const}$, поэтому $p_2=2p_1$. Следовательно, $p_2-p_1=p_1=75\,\text{кПа}$. Ответ: $75\,\text{кПа}$.
}

% Source EE7CA0: https://ege.fipi.ru/bank/questions.php?search=1&proj=BA1F39653304A5B041B656915DC36B38&qid=EE7CA0
\BankTask{EE7CA0}{
Идеальный одноатомный газ находится в сосуде с жёсткими стенками объёмом \(0{,}6\ \text{м}^{3}\). При охлаждении его давление снизилось на \(3\cdot10^{3}\ \text{Па}\).

На сколько уменьшилась внутренняя энергия газа (в $\text{кДж}$)?
}{
\smallskip\noindent\textbf{Решение от GPT}\par
Для одноатомного идеального газа \(U=\frac32pV\). При постоянном объёме уменьшение внутренней энергии равно
\[
U_1-U_2=\frac32V(p_1-p_2)=\frac32\cdot0{,}6\cdot3\cdot10^3=2700\ \text{Дж}=2{,}7\ \text{кДж}.
\]
Ответ: \(2{,}7\ \text{кДж}\).
}

% Source 8054C4: https://ege.fipi.ru/bank/questions.php?search=1&proj=BA1F39653304A5B041B656915DC36B38&qid=8054C4
\BankTask{8054C4}{
Аргон массой $160\,\text{г}$ при комнатной температуре помещают в герметичный закрытый сосуд объёмом $20\,\text{л}$ и начинают нагревать. Масса газа в сосуде остаётся неизменной.

Установите соответствие между физическими величинами, характеризующими аргон, и формулами, выражающими их зависимость от абсолютной температуры $T$ газа в условиях данной задачи.

\begin{center}
\TaskMatchingColumns[.53\linewidth]{Физические величины}{
А) давление газа $p(T)$\par
Б) внутренняя энергия газа $U(T)$
}{Формулы}{
1) $aT$, $a=1662\,\dfrac{\text{Па}}{\text{К}}$\par
2) $bT$, $b=49{,}86\,\dfrac{\text{Дж}}{\text{К}}$\par
3) $\dfrac{c}{T}$, $c=80\,\text{Дж}\cdot\text{К}$\par
4) $dT$, $d=4{,}986\,\dfrac{\text{Па}}{\text{К}}$
}
\end{center}

\begin{center}
\begin{tabular}{|c|c|}\hline
А & Б \\ \hline
\phantom{00} & \phantom{00} \\ \hline
\end{tabular}
\end{center}
}{
\smallskip\noindent\textbf{Решение от GPT}\par
Используем молярную массу аргона $M=40\,\text{г/моль}$. Тогда $\nu=m/M=4\,\text{моль}$, $V=0{,}020\,\text{м}^3$.
\[
p(T)=\frac{\nu R}{V}T=\left(1662\,\frac{\text{Па}}{\text{К}}\right)T,
\qquad
U(T)=\frac32\nu RT=\left(49{,}86\,\frac{\text{Дж}}{\text{К}}\right)T.
\]
Давлению соответствует формула 1, внутренней энергии одноатомного газа — формула 2.

Ответ: А — 1, Б — 2.
}

% Source 1D2B0A: https://ege.fipi.ru/bank/questions.php?search=1&proj=BA1F39653304A5B041B656915DC36B38&qid=1D2B0A
\BankTask{1D2B0A}{
Постоянная масса одноатомного идеального газа в изохорном процессе отдаёт количество теплоты \(Q>0\). Как меняются в этом процессе объём и внутренняя энергия газа?

Для каждой величины определите соответствующий характер изменения:

\begin{enumerate}
\item увеличивается
\item уменьшается
\item не изменяется
\end{enumerate}


\TaskChangeTable{Объём газа}{Внутренняя энергия газа}
}{
\smallskip\noindent\textbf{Решение от GPT}\par
В изохорном процессе объём не изменяется (\(3\)), работа газа \(A=0\). Газ отдаёт теплоту, поэтому полученная им теплота равна \(-Q\). По первому закону термодинамики \(\Delta U=-Q-A=-Q<0\): внутренняя энергия уменьшается (\(2\)).

Ответ: \(32\).
}

% Source 2986: https://phys-ege.sdamgia.ru/problem?id=2986
\BankTask{2986}{
В горизонтальной трубке постоянного сечения, запаянной с одного конца, помещен столбик ртути длиной $15\,\text{см}$, который отделяет воздух в трубке от атмосферы. Трубку расположили вертикально запаянным концом вниз и нагрели на $60$ К. При этом объем, занимаемый воздухом, не изменился. Давление атмосферы в лаборатории  — $750\,\text{мм}$ рт. ст. Какова температура воздуха в лаборатории?
}{
\smallskip\noindent\textbf{Решение}\par
На столбик ртути действуют сила тяжести столбика ртути $mg=\rho gdS,$ сила давления атмосферы $p_0=\rho gHS,$ где $\,\text{Н} = 750\,\text{мм},$ S  — площадь сечения трубки, сила давления воздуха внутри трубки $pS.$
Условие механического равновесия столбика ртути  в вертикальной трубке: 

\[p=p_0 + \rho gd.\]

Поскольку нагревание воздуха в трубке происходит до температуры $T=T_0 + \Delta T$ при неизменном объеме, то по уравнению Клапейрона  — Менделеева: 

\[\frac{T}{T_0} = \frac{p}{p_0} = \frac{\rho gH + \rho gd}{\rho gH} =1 + \frac{d}{H} .\]

Отсюда $T_0=\Delta T \frac{H}{d} = \frac{60 \cdot 750}{150} =300\,\text{К} .$ 

Ответ: $T_0=300K.$
}

% Source 25414: https://phys-ege.sdamgia.ru/problem?id=25414
\BankTask{25414}{
В закрытом сосуде при температуре $100 ^\circ$ находится влажный воздух с относительной влажностью воздуха $60$\% под давлением $100$ кПа. Объем сосуда изотермически уменьшили в $2{,}5$ раза. Во сколько раз нужно увеличить абсолютную температуру воздуха без изменения объема сосуда, чтобы получить такое же конечное давление? Объемом сконденсировавшейся воды пренебречь.
}{
\smallskip\noindent\textbf{Решение}\par
Влажный воздух представляет собой смесь воздуха и водяных паров. Так как температура воздуха $100 ^\circ\,\text{С} ,$ то давление насыщенного водяного пара равно $100$ кПа, в данном случае  — давлению влажного воздуха р.
По закону Дальтона давление газа в сосуде равно сумме парциальных давлений: $p=p_1\,\text{возд} + p_{\,\text{пара}}.$
Из формулы относительной влажности воздуха следует, что давление пара $p_{\,\text{пара}}=\varphi p_{\,\text{н}}. \,\text{п}..$ Откуда давление сухого воздуха в сосуде $p_1\,\text{возд}=p ( 1 - \varphi ) .$
После изотермического сжатия водяной пар стал насыщенным, его давление стало равным давлению насыщенного пара. Давление сухого воздуха при изотермическом сжатии в $2{,}5$ раза увеличилось в $2{,}5$ раза по закону Бойля-Мариотта $pV=const,$ $p_2\,\text{возд}=2{,}5p_1\,\text{возд}.$ Тогда давление влажного воздуха в сосуде стало $p_2=p_2\,\text{возд} + p_{\,\text{н}}. \,\text{п}..$ Таким образом, $p_2=2{,}5 ( 1 - \varphi ) p + p= ( 3{,}5 - 2{,}5\varphi ) p.$
При изохорном нагревании по закону Шарля $\frac{p}{T} =const.$ Откуда:

\[\frac{T_2}{T_1} = \frac{p_2}{p} =3{,}5 - 2{,}5\varphi=3{,}5 - 2{,}5 \cdot 0{,}6=2.\]

Значит, абсолютную температуру воздуха нужно увеличить в $2$ раза.

Ответ: в $2$ раза.
}

\end{document}
